% 编译提示：使用 XeLaTeX 编译
% 需要 ctex 宏包支持中文

\documentclass[a4paper,9pt]{article}
\usepackage{amsmath}
\usepackage{amssymb}
\usepackage{geometry}
\usepackage{ctex}
\usepackage{xcolor}
\usepackage{multicol}
\usepackage{titlesec}
\usepackage{fancyhdr}
\usepackage{tcolorbox}
\usepackage{graphicx}
\usepackage[version=4]{mhchem}

\geometry{left=1cm,right=1cm,top=1cm,bottom=1.8cm}
\setlength{\columnsep}{0.8cm}
\setlength{\columnseprule}{0.4pt}
\renewcommand{\columnseprulecolor}{\color{gray!50}}

\definecolor{sectioncolor}{RGB}{0,102,204}
\definecolor{subsectioncolor}{RGB}{0,153,153}
\definecolor{keycolor}{RGB}{204,102,0}
\definecolor{derivcolor}{RGB}{100,149,237}
\definecolor{boxbg}{RGB}{240,248,255}
\definecolor{keybg}{RGB}{255,248,240}

\titleformat{\section}{\normalfont\large\bfseries\color{sectioncolor}}{\thesection}{0.5em}{}
\titlespacing*{\section}{0pt}{1ex plus 0.5ex minus 0.2ex}{0.8ex plus 0.2ex}
\titleformat{\subsection}{\normalfont\normalsize\bfseries\color{subsectioncolor}}{\thesubsection}{0.5em}{}
\titlespacing*{\subsection}{0pt}{0.8ex plus 0.3ex minus 0.1ex}{0.5ex plus 0.1ex}
\setcounter{secnumdepth}{4}\setcounter{tocdepth}{4}

\newenvironment{keypoint}
{\par\vspace{0.3em}\noindent\begin{tcolorbox}[colback=keybg,colframe=keycolor,boxrule=0.5pt,arc=2pt,left=2pt,right=2pt,top=2pt,bottom=2pt,boxsep=0pt]\noindent\textcolor{keycolor}{\textbf{关键点：}}}
{\end{tcolorbox}\par\vspace{0.3em}}

\newenvironment{examplebox}
{\par\vspace{0.3em}\noindent\begin{tcolorbox}[colback=boxbg,colframe=derivcolor,boxrule=0.5pt,arc=2pt,left=2pt,right=2pt,top=2pt,bottom=2pt,boxsep=0pt]\scriptsize\noindent\textcolor{derivcolor}{\textbf{示例：}}\par\setlength{\parskip}{0.1ex}\setlength{\itemsep}{0pt}\setlength{\parsep}{0pt}}
{\end{tcolorbox}\par\vspace{0.3em}}

\pagestyle{fancy}
\fancyhf{}
\fancyfoot[C]{\scriptsize9-\thepage}
\renewcommand{\headrulewidth}{0pt}
\renewcommand{\footrulewidth}{0.4pt}
\renewcommand{\footrule}{\hbox to\headwidth{\color{gray!50}\leaders\hrule height \footrulewidth\hfill}}

\title{\vspace{-2em}\Large\bfseries\color{sectioncolor} Chapter 9 Collections - 表面吸附、势能面与过渡态（Lect.10）\vspace{-1em}}
\author{}
\date{\today}

\begin{document}
\maketitle
\thispagestyle{fancy}

\begin{multicols}{2}
\scriptsize{\tableofcontents}

\section{气体在表面的吸附}

\subsection{Langmuir 模型假设}
设表面共有 $M$ 个\textbf{等价吸附位点}，其中有 $N$ 个被占据，覆盖度定义为
\[
\theta=\frac{N}{M}.
\]
Langmuir 模型的基本假设是：
\begin{itemize}
  \item 每个位点至多吸附 1 个分子；
  \item 各位点彼此等价；
  \item 分子是否吸附在某位点，不受邻位占据情况影响。
\end{itemize}

\begin{keypoint}
这是一个\textbf{单层吸附 + 位点独立}模型，因此其成功与否取决于真实表面是否足够接近“均匀、无相互作用位点”的理想情况。
\end{keypoint}

\subsection{Langmuir 等温式}
在该模型下，表面覆盖度满足
\[
\theta=\frac{K_{\mathrm{ads}}p}{1+K_{\mathrm{ads}}p}.
\]
低压极限 $K_{\mathrm{ads}}p\ll 1$ 时，
\[
\theta\approx K_{\mathrm{ads}}p,
\]
高压极限 $K_{\mathrm{ads}}p\gg 1$ 时，
\[
\theta\to 1.
\]
其线性化形式为
\[
\frac{p}{\theta}=p+\frac{1}{K_{\mathrm{ads}}},
\]
因此作图 $p/\theta$ 对 $p$ 应为直线，斜率为 1，截距为 $1/K_{\mathrm{ads}}$。

\begin{examplebox}
Ex.~35 的动力学推导：设吸附速率为
\[
r_{\mathrm{ads}}=k_a(M-N)p,
\]
脱附速率为
\[
r_{\mathrm{des}}=k_dN.
\]
平衡时二者相等，得
\[
K_{\mathrm{ads}}=\frac{k_a}{k_d}=\frac{\theta}{(1-\theta)p},
\]
于是直接整理成 Langmuir 等温式；若题目要求实验判据，则继续化成
\[
\frac{p}{\theta}=p+\frac{1}{K_{\mathrm{ads}}}.
\]
\end{examplebox}

\section{Langmuir 式的统计热力学推导}

\subsection{从排列数到系统配分函数}
把 $N$ 个\textbf{不可分辨}分子放到 $M$ 个\textbf{可区分}位点上的方法数是
\[
\Omega=\frac{M!}{N!(M-N)!}.
\]
若单个吸附分子的内部配分函数为 $q$，则整体系统配分函数可写成
\[
Q_N=q^N\frac{M!}{N!(M-N)!}.
\]
于是 Helmholtz 自由能为
\[
A=-kT\ln Q_N.
\]

对阶乘用 Stirling 近似，得到
\[
A=kT\left[-N\ln q-M\ln M+N\ln N+(M-N)\ln(M-N)\right].
\]

\subsection{吸附态化学势}
吸附相化学势由
\[
\mu_{\mathrm{ads}}=\left(\frac{\partial A}{\partial N}\right)_{T,M}
\]
给出，因此
\[
\mu_{\mathrm{ads}}=-kT\ln q+kT\ln\left(\frac{N}{M-N}\right)
=-kT\ln q+kT\ln\left(\frac{\theta}{1-\theta}\right).
\]

而气相分子的化学势写成
\[
\mu_g=\mu^\circ+kT\ln\frac{p}{p^\circ}.
\]
平衡条件是
\[
\mu_{\mathrm{ads}}=\mu_g.
\]
整理后可得
\[
b\frac{p}{p^\circ}=\frac{\theta}{1-\theta},
\qquad
b=q\exp\left(\frac{\mu^\circ}{kT}\right),
\]
于是
\[
	\theta=\frac{K_{\mathrm{ads}}p}{1+K_{\mathrm{ads}}p}.
\]

\begin{keypoint}
	这一段有两套等价但量纲处理不同的写法：
	\begin{itemize}
		\item 若从 $\mu_g=\mu^\circ+kT\ln(p/p^\circ)$ 出发，则对数里的 $p/p^\circ$ 必须无量纲，此时 $b=q\exp(\mu^\circ/kT)$ 也是无量纲；
		\item 若改写成
		\[
			\theta=\frac{K_{\mathrm{ads}}p}{1+K_{\mathrm{ads}}p},
		\]
		则 $K_{\mathrm{ads}}$ 必须带 $1/\text{pressure}$ 单位（如 $\mathrm{bar^{-1}}$、$\mathrm{Pa^{-1}}$）。
	\end{itemize}
	因此做题时不要把“无量纲的 $b$”和“有压强倒数量纲的 $K_{\mathrm{ads}}$”混成同一个常数。
\end{keypoint}

\begin{keypoint}
	Ex.~36 的本质是：\textbf{位点组合数}决定 $Q_N$，$Q_N$ 决定 $A$，$A$ 对 $N$ 的偏导给出吸附相化学势，最后由化学势平衡导出等温式。
\end{keypoint}

\section{速率常数的概念铺垫}

统计热力学既然能从微观能级推平衡常数，下一步自然会问：能否同样处理\textbf{速率常数}？lect10 的回答是：可以，但必须先引入\textbf{势能面}和\textbf{过渡态}。

\section{势能面（PES）}

\subsection{PES 是什么}
势能面是电子能量对核坐标的函数。对任一原子排布，只要给定核位置，就可由量子化学原理上求得其电子能量。因此反应从反应物到产物的全过程，都对应 PES 上的一条路径。

PES 上的\textbf{极小值}对应稳定分子：当几何略微扰动时，势能升高，系统有恢复力把分子拉回平衡结构。

\subsection{$A+BC\to AB+C$ 的简化图像}
若把最简单反应限制为线型，体系只需两个坐标：
\[
r_{AB},\qquad r_{BC}.
\]
于是 PES 可以画成两变量函数。反应物位于一条谷中，产物位于另一条谷中；两条谷在中部相连，形成一个“山口”。

\begin{keypoint}
对真实多原子体系，反应坐标通常是许多核运动的复合，不是某一根键长单独变化；但这个二维模型足够建立考试中最重要的图像。
\end{keypoint}

\section{反应路径、过渡态与活化能}

\subsection{最小能量路径}
在 PES 上从反应物到产物可以走很多路线，但绝大多数成功反应都沿着\textbf{最低山口}通过，也就是\textbf{最小能量路径}（minimum energy pathway）。因为若某条路径需要更高能量，能达到它的分子只会更少。

\subsection{过渡态与中间体}
最小能量路径上的最高点就是\textbf{过渡态}（transition state, TS）。它满足：
\begin{itemize}
  \item 沿反应坐标方向是势能极大值；
  \item 在其余方向上是极小值。
\end{itemize}
因此 TS 不是普通稳定分子，而是只有极短寿命的瞬态，大约只有一个分子振动周期量级（$\sim 10^{-13}\,\mathrm s$）。

这与\textbf{中间体}不同：中间体位于 PES 极小值，是真正稳定到可定义、可探测的分子种类。

\begin{keypoint}
一句话区分：\textbf{中间体在谷底，过渡态在山口。}
中间体可存在，过渡态只能“经过”。
\end{keypoint}

\subsection{反应坐标与活化能}
最小能量路径本身就叫\textbf{反应坐标}。若把能量沿这条路径投影出来，就得到熟悉的反应剖面图。过渡态相对反应物的能量差定义为活化能：
\[
E_a=E_{\mathrm{TS}}-E_{\mathrm{reactants}}.
\]

\begin{examplebox}
PES 部分最常见的题不是计算，而是概念判断：
\begin{itemize}
  \item 哪个点是 reactants / products / TS / intermediate；
  \item 为什么反应走 minimum-energy path；
  \item 为什么 TS 沿一方向不稳定、而其余方向稳定；
  \item 如何从反应剖面图上读出 $E_a$。
\end{itemize}
\end{examplebox}

\section{最后速记}

\begin{keypoint}
lect10 高频公式与结论：
\[
\theta=\frac{K_{\mathrm{ads}}p}{1+K_{\mathrm{ads}}p},
\qquad
\theta\approx K_{\mathrm{ads}}p\ (\text{低压}),
\qquad
\theta\to 1\ (\text{高压}),
\]
\[
\begin{aligned}
Q_N&=q^N\frac{M!}{N!(M-N)!},\\
A&=-kT\ln Q_N,\\
\mu_{\mathrm{ads}}&=-kT\ln q+kT\ln\frac{\theta}{1-\theta}.
\end{aligned}
\]
\[
\mu_{\mathrm{ads}}=\mu_g,
\qquad
\mu_g=\mu^\circ+kT\ln\frac{p}{p^\circ},
\]
\[
\text{TS = 最小能量路径最高点},
\qquad
E_a=E_{\mathrm{TS}}-E_{\mathrm{reactants}}.
\]
\end{keypoint}

\end{multicols}
\end{document}
